\section{Caste and Spirit}

Caste refers to the division of society into functional and hierarchical classes reflecting the different characteristics of human individuals. ``Caste'' is often exclusively associated with the societal divisions of Hindu India, although the word originally referred to the feudal system in Portugal. Furthermore, George Dumezil documented the caste structure of Traditional Indo-European societies from India to Ireland.

In keeping with Guenon's and Evola's usage, we will usually refer to the Hindu names of the castes. The Brahmin caste refers to those who are oriented to transcendent reality; their function is to maintain spiritual authority and ensure the conservation and transmission of traditional doctrine. The Kshatriyas consist of the ruling and warrior class. The Vaishyas comprise the productive elements of society. Finally, there is the mass, or Shudra caste, and beyond them the outcastes, which may be slaves, foreigners, or others inhabiting the land, yet outside the societal bounds. This system was last seen in the West during the Middle Ages where the clergy corresponded to the Brahmins, the nobility to the Kshatriyas, the third-estate to the Vaishyas, and the serfs to the Shudras.

Evola ties the notion of castes to aryanity, or the capacity to be twice born.

\begin{quotex}
In a specific sense, ``Aryan'' was essentially a designation of caste: it referred collectively to the whole of the three higher castes (spiritual leaders, warrior aristocracy, and ``heads of families'' or landowners with authority over a certain group related by blood) in their opposition to the fourth caste of the Shudra -- today, perhaps, we would say the proletarian masses.

\flright{\textsc{Julius Evola}, \textit{Sintesi di dottrina della razza}}

\end{quotex}

Although the Brahmins are the highest caste, they do not rule by power, but rather by moral authority. They enunciate the transcendent principles that rule the society. The Kshatriyas have the political power. Their role is to apply the principles to specific circumstances. They derive their authority from the Brahmins; in return, they provide a safe and stable milieu for the Brahmins to do their own work.

Due to their inherent differences, the different castes experience spirituality in different ways; the system or religion needs to be broad enough to accommodate the different perspectives. The Brahmins are oriented toward super-human states and apply themselves, in their personal practice, to the realization of final Deliverance.

The path of the Kshatriyas is more devotional or bhakti, that is, it is more emotional than the intellectual path of the Brahmin. The Kshatriya is interested in developing the highest potentials of the human state. The Vaishya is interested in moral stability, the preservation of property, and the maintenance of family and clan loyalties. The Shudras believe anything; however, whenever the caste system is weakened, they turn to atheism, since hierarchical structure is a reflection of divine order.

In a Traditional, hierarchically organized society, the religious form is regulated by the true spiritual authority, or Brahmins. Since the Medieval period, there have been a succession of revolts against this authority by each lower caste, until we have reached our current state which is ruled by functional atheism and materialism.


\flrightit{Posted on 2010-03-03 by Cologero}

